{{دیوان
| عنوان = اُمِّید
| ماخذ = ویکی ماخذ
| ماخذ_لنک = https://ur.wikisource.org/wiki/%D8%A7%D9%8F%D9%85%D9%90%D9%91%DB%8C%D8%AF
| مقابلہ = اقبال سائبر لائبریری (اقبال اکادمی پاکستان)
| مقابلہ_لنک = https://www.iqbalcyberlibrary.net
}}
<poem>
مقابلہ تو زمانے کا خوب کرتا ہُوں
اگرچہ مَیں نہ سپاہی ہُوں نے امیرِ جنُود

مجھے خبر نہیں یہ شاعری ہے یا کچھ اور
عطا ہُوا ہے مجھے ذکر و فکر و جذب و سرود

جبینِ بندۂ حق میں نمود ہے جس کی
اُسی جلال سے لبریز ہے ضمیرِ وجود

یہ کافری تو نہیں، کافری سے کم بھی نہیں
کہ مردِ حق ہو گرفتارِ حاضر و موجود

غمیں نہ ہو کہ بہت دَور ہیں ابھی باقی
نئے ستاروں سے خالی نہیں سِپہرِ کبود
</poem>

٭ ریاض منزل (دولت کدۂ سرراس مسعود) بھوپال میں لکھے گئے
